% ============================================================================
% SEC05.TEX - Section 8.5: Menu System
% ============================================================================

\section{Menu System}

\begin{learningoutcomes}
\item Membuat Main Menu
\item Mengimplementasikan Pause Menu
\item Menggunakan SceneManager untuk navigation
\end{learningoutcomes}

\subsection{Main Menu}

\begin{lstlisting}[language={[Sharp]C}]
using UnityEngine;
using UnityEngine.SceneManagement;
using UnityEngine.UI;

public class MainMenu : MonoBehaviour
{
    public Button startButton;
    public Button settingsButton;
    public Button quitButton;
    
    void Start()
    {
        startButton.onClick.AddListener(StartGame);
        settingsButton.onClick.AddListener(OpenSettings);
        quitButton.onClick.AddListener(QuitGame);
    }
    
    void StartGame()
    {
        SceneManager.LoadScene("GameScene");
    }
    
    void OpenSettings()
    {
        // Open settings menu
    }
    
    void QuitGame()
    {
        Application.Quit();
    }
}
\end{lstlisting}

\subsection{Pause Menu}

\begin{lstlisting}[language={[Sharp]C}]
public class PauseMenu : MonoBehaviour
{
    public GameObject pausePanel;
    private bool isPaused = false;
    
    void Update()
    {
        if (Input.GetKeyDown(KeyCode.Escape))
        {
            TogglePause();
        }
    }
    
    void TogglePause()
    {
        isPaused = !isPaused;
        pausePanel.SetActive(isPaused);
        Time.timeScale = isPaused ? 0f : 1f;
    }
}
\end{lstlisting}

\begin{catatan}
Gunakan Time.timeScale = 0 untuk pause game. Pastikan untuk reset ke 1 saat resume. Gunakan separate scene atau Canvas untuk menu.
\end{catatan}

\subsection{Ringkasan Section}

\begin{itemize}
\item Main Menu untuk memulai game
\item Pause Menu untuk pause/resume gameplay
\item Time.timeScale untuk pause functionality
\item SceneManager untuk navigation antar scene
\end{itemize}
